\documentclass[10pt]{article}
\usepackage[margin=20mm]{geometry}
\usepackage{amsmath,amssymb,mathtools,fontspec,unicode-math,xurl,hyperref}
\IfFontExistsTF{DejaVu Serif}{\setmainfont{DejaVu Serif}}{\setmainfont{Latin Modern Roman}}
\setmathfont{latinmodern-math.otf}
\setlength{\parindent}{0pt}\setlength{\parskip}{6pt}\sloppy
\title{Initial technical review of
Cluster-size decay in supercritical
kernel-based spatial random graphs}
\author{AI-assisted technical review}
\date{}
\begin{document}
\maketitle
Abstract: Theorem 2.4 is false on part of its stated lattice sum-kernel range: the written parameter restriction permits a deterministic connected nearest-neighbor subgraph, whereas the theorem assigns positive probability to an impossible lower-tail event at arbitrarily large volumes. This is a Major range defect. Eight Minor problems concern an introductory asymptotic sign, finite-volume endpoints, a degenerate boundary exponent, a dropped mark restriction, a Palm comparison, a conditioning field, an auxiliary monotonicity claim, and a missing cover-expansion case. Explicit evidence and repairs are given. The identified proof defects do not establish failure of the Poisson cluster-tail or second-component asymptotics. The review also checks the main exponent algebra, coupling, isolation mechanism and cited hypotheses, with version and verification limits stated separately.\par
\section*{Reviewed object, hypotheses and scope}
The assigned object is M25-0518, by Joost Jorritsma, Júlia Komjáthy and Dieter Mitsche, DOI 10.1214/24-AOP1742 (https://doi.org/10.1214/24-AOP1742) [JKM]. The exact acquired PDF is arXiv:2303.00724v3, 17 October 2024, with 74 physical pages. Physical and printed page numbers agree. All source locators in this report refer to this version. Its SHA256 is\par
> `dcf5cdc87d169a282eeeea4b7daa8ab5b8a3aceb5526b2caa1f851920bba44f8`.\par
The scientific decision remains pending.\par
Equation reference: model. The vertex locations form either a unit-intensity Poisson process in \(\mathbb R^d\) or the lattice \(\mathbb Z^d\). Marks are independent, with \(\mathbb P(W>w)=w^{-(\tau-1)}\) for \(w\geq1,\ \tau>2\), or \(W\equiv1\) when \(\tau=\infty\). Conditional on locations and marks, distinct edges are independent and have probability \[
 q(u,v)=p\left(1\wedge
 \frac{\beta\kappa(w_u,w_v)}{\|x_u-x_v\|^d}\right)^\alpha,
 \qquad \alpha>1,\quad p\in(0,1],\quad\beta>0.\] For \(\alpha=\infty\), the power profile is replaced by its threshold version. The kernels are \(\kappa_\sigma(w,w')=(w\vee w')(w\wedge w')^\sigma\), \(\sigma\geq0\), and \(\kappa_{\rm sum}(w,w')=(w^{1/d}+{w'}^{1/d})^d\). Assumption 1.3 adds \(p\wedge\beta<1\) on the lattice. The intended box \(\Lambda_n=[-n^{1/d}/2,n^{1/d}/2]^d\) has volume \(n\), which is not the exact number of lattice vertices.\par
Write \(t=\tau-1\). The four candidate boundary exponents are \[\zeta_{\rm short}=(d-1)/d,\qquad
 \zeta_{\rm ll}=2-\alpha,\qquad
 \zeta_{\rm hl}=t/\alpha-(t-1),\qquad
 \zeta_{\rm hh}=1-t\gamma_{\rm hh}.\] For finite \(\alpha\) and \(\tau\leq\sigma+2\), \(\gamma_{\rm hh}=(1-1/\alpha)/(\sigma+1-t/\alpha)\); the other branch is \(1/(\sigma+1)\). The source defines \(\zeta_\star\) as their maximum and \(m_\star\) as the number of maximizing entries. The main matching cluster-tail and second-component bounds are developed when \(\zeta_{\rm hh}>0\), equivalently \(2<\tau<\sigma+2\). Their upper bounds are for Poisson locations, with a separate exponent when \(\tau<\sigma+1\). The lower-tail Theorem 2.4 and Claim 1.4 have a wider stated range. That distinction is essential to the findings below.\par
All 74 source pages, including Appendices A–C and the references, were read. Actual source-image inspection covered 58 pages, including every page used as evidence for the numbered problems and the main counterevidence. Twenty pages of three acquired references were also read and viewed. Reading an index or parsing a manifest is not counted as mathematical reading. The review checks the finite-volume statements, the cover construction, high-mark backbone and conditional coupling, the cluster-tail reduction, the isolation lower bound, and the appendix calculations. It is not a formal proof certification of every external theorem or every extension left to the reader.\par
\section*{A false theorem range and endpoint assertions}
Problem 1 (Major: the lattice sum-kernel restriction permits certain edges).  Theorem 2.4, p. 10, (2.8), is false under the literal lattice hypotheses of Assumption 1.3, p. 4. The relevant proof uses the strict edge-probability assertion in the lattice adaptation of Lemma 7.2, p. 49, and then the covering argument on pp. 52–53. The comparison \(\kappa_{\rm sum}\leq2^d\kappa_0\) on p. 12 does not preserve the stated nondegeneracy restriction.\par
Equation reference: zeroevent. Evidence. Take \(d=2,\ W\equiv1,\ \kappa=\kappa_{\rm sum},\
\alpha=2,\ p=1,\ \beta=3/4\). These choices satisfy the displayed assumptions, including \(p\wedge\beta<1\). But \(\kappa_{\rm sum}(1,1)=4\), so every lattice nearest-neighbor edge has probability \((1\wedge3)^2=1\). For \(n=4m^2\), \(m\geq1\), the box is \([-m,m]^2\). Its \((2m+1)^2\) lattice vertices form a connected deterministic nearest-neighbor graph. Consequently, for \(\rho=1/2\), \[
 \mathbb P(|C_n^{(1)}|<\rho n)=0.\] Here \(\zeta_\star=1/2\) and \(m_\star=1\): the short-edge exponent dominates \(2-\alpha=0\), and the heavy-mark mechanisms disappear for unit marks. The right side of (2.8) is \(\exp[-(A/\rho)n^{1/2}]=\exp(-4Am)>0\) for every finite \(A\). Thus the contradiction holds at arbitrarily large volumes. ◻\par
The phrase explaining that the graph should not be almost surely connected can plausibly express an intended extra hypothesis. It does not follow from the written restriction for this kernel. Suppressing marks near their minimum cannot remove certain nearest-neighbor edges. Moreover, comparison with \(\kappa_0\) changes the effective parameter to \(2^d\beta=3\), precisely outside the claimed lattice restriction when \(p=1\). A sufficient local correction is to require \(p<1\) or \(\beta\kappa(1,1)<1\), with the close-to-one mark cutoff in the isolation proof chosen accordingly. For \(\kappa_{\rm sum}\), this means \(p<1\) or \(2^d\beta<1\). The correction must be propagated through the lattice statement and proof. This is a false final assertion on a stated subfamily, not merely a missing estimate. It does not refute the Poisson upper bounds or the heavy-heavy kernel results on their nondegenerate range.\par
Problem 2 (Minor: the introductory lower-tail order has the wrong direction).  Theorem 1.1, p. 2, (1.3), gives \(\log\mathbb P(|C_n^{(1)}|<\rho n)=-\Omega(n^\zeta)\) for every \(\rho>0\). The detailed Theorem 2.4 gives a lower bound on this probability, and the discussion on p. 11 and the proof of Theorem 1.1 on p. 54 identify that lower bound as the intended input.\par
For an explicit permitted heavy-heavy example, take Poisson locations, \(d=1,\sigma=1,\tau=5/2,\alpha=2\), and any fixed \(p\in(0,1]\), \(\beta>0\). Then \(\zeta_{\rm hh}=2/5>\max(\zeta_{\rm ll},\zeta_{\rm short})=0\), as required in the introductory statement. Let \(\rho=2\). Since the number of vertices \(N_n\) is \({\rm Poi}(n)\), \[\mathbb P(|C_n^{(1)}|<2n)\geq \mathbb P(N_n<2n)
 \geq1-\exp[-(2\log2-1)n].\] The logarithm tends to zero and cannot be bounded above by \(-c n^{2/5}\), \(c>0\). The lower bound actually proved in Theorem 2.4 implies \(\log\mathbb P(\cdot)=-O(n^\zeta)\) for each fixed \(\rho\) in the unique-maximizer setting of Theorem 1.1. Replacing \(\Omega\) by \(O\) is consistent with both the proof and the prose. An upper large-deviation estimate would instead need a restriction such as \(\rho<\theta\) and its own argument. This correction concerns the introductory formulation, not the cluster-tail exponent.\par
Problem 3 (Minor: the all-size quantifiers fail at logarithmic and lattice endpoints).  Theorem 2.1(i), p. 9, states its lower bound for \(k\geq1\); Theorem 2.4, p. 10, states (2.8) for \(n\geq1\) and \(\rho>0\). These quantifiers admit impossible lower bounds.\par
Take \(d=1,\sigma=1,\tau=5/2,\alpha=3/2\). Then \(\zeta_{\rm ll}=\zeta_{\rm hl}=\zeta_{\rm hh}=1/2\), \(\zeta_{\rm short}=0\), and \(m_\star=3\). For Poisson locations and \(k=1\), the claimed right side of the cluster lower bound is \(\exp[-A\,1^{1/2}(\log1)^2]=1\). At any finite \(n>\max(A,1)\), however, the Palm origin is the sole vertex in the box with probability \(e^{-n}>0\). On that event \(|C_n(0)|>1\) fails, irrespective of the additional condition that the origin be outside the largest component. Thus the left side is strictly below one.\par
For Theorem 2.4 with the same parameters, \(n=\rho=1\), the right side is again one, while \(\mathbb P(|C_1^{(1)}|<1)=\mathbb P(N_1=0)=e^{-1}\). Independently of logarithmic ties, the lattice always contains the origin; if \(\rho n\leq1\), its lower-tail event is empty, whereas a finite exponential right side is positive. The concentration arguments on pp. 50 and 53 are naturally large-size arguments. Stating the cluster estimate for sufficiently large \(k\), and the giant estimate for sufficiently large \(n\) depending on \(\rho\), removes these particular endpoint contradictions after the nondegeneracy correction in Problem Problem 1 (\#prob:lattice). A replacement such as \(\log(2+k)\) addresses the logarithmic endpoint alone; it cannot make an impossible lattice event have positive probability. No asymptotic cluster-tail conclusion is disproved by this problem.\par
Problem 4 (Minor: the boundary exponent needs a nonempty-edge qualification).  Claim 1.4, p. 7, and Lemma 7.7, p. 47, (7.27), include the threshold and unit-mark limits without a separate nondegeneracy condition. The short-edge lower-bound sketch on p. 69 assumes that a boundary vertex has a positive chance of an edge across the boundary.\par
Take \(d=2,\ W\equiv1,\ \kappa=\kappa_0=1,\
\alpha=\infty,\ p=1,\beta=3/4\), on \(\mathbb Z^2\). Every distinct pair has \(\|x-y\|^d\geq1>\beta\); hence there are no edges. Assumption 1.3 still holds. The expected downward vertex-boundary size is zero, so its logarithm is undefined as a real number (or is \(-\infty\) in the extended convention). The displayed maximum predicts \(\zeta_{\rm short}=1/2\). This is a separate degenerate case from the connected example in Problem Problem 1 (\#prob:lattice). Restrict the claim to a setting with a positive crossing-edge mechanism, or state the empty case separately. Claim 1.4 is not stated only in the supercritical heavy-heavy range; that narrower context would exclude this example. The main supercritical asymptotic theorem is therefore not refuted by it.\par
\section*{Proof dependencies requiring local repairs}
Problem 5 (Minor: a component-wide mark restriction is dropped).  The definition of \(A_{\rm low\text{-}edge}\) in (6.10), p. 36, bounds the two edge-endpoint marks. The event inclusion in the proof of Claim 6.5, p. 39 immediately before (6.24), implicitly needs all vertices on the connecting path to have low marks. This stronger condition is absent from that event.\par
Evidence and repair. Fix the finite-\(\alpha\) case, with \(p\in(0,1)\), and take \(k\) sufficiently large. Write \(t_k\) for the length cutoff and \(\bar w\) for the mark cutoff. In dimension one put three vertices near \(0,2t_k,4t_k\), inside the box of half-side \(k t_k\). Give the origin mark greater than \(\bar w\), and give the other two marks less than \(\bar w\). Let the only component edges be \(0a\) and \(ab\). Both have length greater than \(t_k\). The full component has size three, at most \(k\), and the low-low edge \(ab\) makes \(A_{\rm low\text{-}edge}\) hold. After all edges of length at least \(t_k\) are deleted, the origin’s component is just the high-mark origin. It contains no low-mark endpoint of the kind required on the right side of the displayed inclusion.\par
This is a genuine finite-configuration obstruction: use small open position and mark intervals and the specified finite edge pattern. Under finite \(\alpha\) and \(0<p<1\), that pattern has positive conditional probability. For large \(t_k\) the three points do not violate the annular density bounds; further points may be placed sparsely without changing the component. The counterexample refutes the event implication, not the numerical probability bound (6.22).\par
There is strong counterevidence to a failure of the needed cluster-tail estimate. In (6.19), p. 37, the intersection with \(A_{\rm mark\text{-}giant}\) and the condition of being outside the giant already ensures that every vertex of the relevant component has mark at most \(\bar w\). Retain this restriction when defining the event to be bounded. Along a path from the origin, the first edge of length at least \(t_k\) then has a low-mark endpoint in the short-edge component; the argument before (6.24) applies. Conditional independence of the unrevealed long-edge variables is unchanged. ◻\par
The defect lies in a core upper-bound reduction, but the required stronger event is available immediately upstream. The repair changes the intermediate statement and its use, without supplying a counterexample to Theorem 2.1(ii)–(iii).\par
Problem 6 (Minor: the Palm-to-ordinary comparison needs a degree argument).  On p. 50, between (7.37) and (7.38), the proof replaces the probability of a large ordinary largest component by the Palm probability of a large origin component, losing only one vertex in the threshold. Adding a Palm vertex can merge several ordinary components; deleting it need not leave all but one of its component in a single component.\par
For a concrete obstruction to that general comparison, take a one-dimensional Poisson KSRG in \([-2,2]\), with \(\alpha=\infty,\beta=2,p=1,\sigma=1,\tau=5/2\), without mark truncation. The Palm origin is connected to every ordinary vertex, since distances to it are at most two and all marks are at least one. Hence \[\mathbb P^0(|C_4(0)|>8)=\mathbb P({\rm Poi}(4)\geq8).\] The ordinary event \(|C_4^{(1)}|>7\) is contained in \(\{N_4\geq8\}\). It misses the positive-probability event with exactly four vertices in each of \([-2,-1.8]\) and \([1.8,2]\), no others, and all eight marks in \([1,1.1]\). The two groups cannot connect: their gap is at least \(3.6\), while \(\beta w w'\leq2(1.1)^2=2.42\). That event has probability \[e^{-4}\frac{(1/5)^8}{(4!)^2}
       \bigl[1-(11/10)^{-3/2}\bigr]^8>0.\] Thus the unqualified one-vertex comparison is false even for a KSRG. This finite-box example does not assert failure of the particular eventual, mark-truncated inequality under the additional hypothesis (7.2). It shows why merely ignoring Palm conditioning is not a proof of that step.\par
A replacement sufficient for the asymptotic lower bound follows from local finiteness. Let \(D_0\) be the degree of the Palm origin in the infinite graph. Under the present kernel, \(\tau>2\) and \(\alpha>1\) imply \(D_0<\infty\) almost surely, also when its unconditional expectation is infinite. Given the positive constant \(\rho\) in (7.2), choose a fixed \(K\) with \(\mathbb P^0(D_0>K)<\rho/2\). On the event that the truncated Palm component has at least \(\rho n\) vertices and \(D_0\leq K\), deletion of the origin leaves at most \(K\) components. One has size at least \((\rho n-1)/K\). The ordinary marked graph has its original law under this deletion. It therefore has a component of a smaller but positive linear density with probability at least \(\rho/2\). Enlarging the outer box by the corresponding fixed factor recovers the outside component needed in (7.7)–(7.38). All exponential powers remain unchanged. The printed constant \(\rho^2\) and threshold cannot simply be retained without this adjustment. Proposition 7.1’s final asymptotic conclusion is not disproved.\par
Problem 7 (Minor: the conditional field in (A.7) reveals the wrong mark bands).  Equation (A.7), p. 56, conditions on \(G_n[1,w_{\rm hh})\). The proof of Claim A.2 instead requires the intermediate mark bands below \(w_{\rm hh}\) to remain random.\par
The relevant hypotheses are Poisson locations, \(\zeta_{\rm hh}>0\), a realized high-mark backbone satisfying \(A_{\rm bb}\), regularity of the lowest band \([m_w,2m_w)\), and the origin mark in that band. For sufficiently large \(n\), choose an intermediate band \(I_j=[2^j m_w,2^{j+1}m_w)\), \(j\geq1\), strictly below \(w_{\rm hh}\). Under the literal conditioning in (A.7), the count in \(Q_j(x)\times I_j\) is already known. A realization with count zero has positive probability and is compatible with the lowest-band regularity and the disjoint high-mark backbone. Its conditional probability of a candidate vertex is zero, whereas the positive-mean Poisson lower domination invoked immediately before (A.17), p. 58, assigns positive probability to at least one candidate. The same conditioning also invalidates the later claim that the revealed low-band positions are still uniform.\par
The last line of p. 57, the explanation on p. 58, and the proofs on pp. 59–60 explicitly describe conditioning only on \(G_n[w_{\rm hh},2w_{\rm hh})\). Replacing the graph in (A.7) by this high-mark graph, while keeping the separately stated lowest-band regularity and origin-mark conditions, restores the independence of the intermediate Poisson bands. This is a locally repairable discrepancy in the definition of a conditional law, not evidence that the linear-component construction is false. In the same argument, the second line of (A.18), p. 59, also needs a minus sign in the exponent: \(1-(1-q)^s\geq1-e^{-sq}\), not \(1-e^{sq}\). The first line supplies the intended Bernoulli bound directly.\par
\paragraph*{Notation clarification: box volume.}\mbox{}\par
Equation (2.10), p. 12, contains a local typographical error in the box notation. It is not counted as a scientific Problem.\par
For \(s>0\) and integer \(d\geq1\), the displayed box \[x+[-s/2^{1/d},s/2^{1/d}]^d\] has volume \(2^{d-1}s^d\), not \(s\). For example \(d=s=2\) gives volume eight. The correct expression is \(x+[-s^{1/d}/2,s^{1/d}/2]^d\). The correctly normalized \(\Lambda_n\) on p. 4 and the explicit box expression on p. 22 are counterevidence to an intended alternative convention. The later tessellation and cover estimates use volume \(s\), so this correction should be made before interpreting those formulas. No altered exponent or false asymptotic theorem follows from this typographical error, so it is excluded from the substantive issue count.\par
Problem 8 (Minor: the auxiliary function \(g\) is not always monotone).  After (B.16), p. 66, the proof asserts that \(g\) is nondecreasing in both arguments and uses this in (B.18). The \(\alpha=\tau-1,\ \sigma\alpha>\tau-1\) row defining \(g_2\) on p. 65 need not have that property.\par
Take \(\alpha=\tau-1=3/2,\sigma=6/5,d=1,\gamma=1/4\) and \(C_\beta=1\), so \(\beta=1/2\). Choose \(r_k\geq100\). In the region \(1\leq u\leq v=100\), the row in question reads \[g(100,u)=\bigl(1+\tfrac14\log(100/u)\bigr)u^{3/40}.\] Its derivative at \(u=100\), from the left, is \[-\frac{7}{40}\,100^{-37/40}
 =-\frac{7\,10^{3/20}}{4000}<0.\] All parameters satisfy the stated assumptions, including \(\gamma\leq1/(\sigma+1)\). The positivity of the power does not make the product with a decreasing logarithmic factor monotone.\par
The estimate is nevertheless repairable with a constant, which is all the \(\lesssim\) notation requires. For \(0<u\leq R\leq v\), \(c,q>0\), put \(h(v,u)=(1+c\log(v/u))u^q\). Then \[\frac{h(v,u)}{h(v,R)}
 \leq (u/R)^q\bigl(1+c\log(R/u)\bigr)
 \leq 1+\frac{c}{e q}.\] The last step uses \(\sup_{0<x\leq1}x^q\log(1/x)=1/(eq)\). For the relevant row \(c=\gamma d\) and \(q=\gamma d(\sigma-1)(\tau-1)>0\); their ratio is fixed by the kernel parameters. This comparison justifies the replacement by \(r_k\) in (B.18) up to a fixed constant, without the false monotonicity claim. It preserves the integral exponents in Lemma 7.4 and the isolation bound.\par
Problem 9 (Minor: Lemma 5.9 omits the \(s(w)>n\) cover case).  Proposition 4.2, p. 15, requires \(s(w)\leq n\), where \(s(w)=(2^d\beta w)^{1/(1-1/\alpha)}\). Its invocation on p. 31 in the proof of Lemma 5.9 checks the lower bound on \(w=2w_{\rm hh}(k)\), but not this upper bound on \(s(w)\). The ambient size may be any \(n\geq k\).\par
For finite \(\alpha=11/10,\tau=5/2,\sigma=1\), \(\gamma_{\rm hh}=1/7\), and therefore \[s(2w_{\rm hh}(k))\asymp k^{11/7}.\] At \(n=k\), the required inequality fails for all sufficiently large \(k\). The failure is a domain-of-application omission, not a counterexample to the merging estimate.\par
Here is the missing case. If \(s(w)>n\), select exactly \(k\) vertices from the component of size greater than \(k\) (the construction uses integer \(k\)). This subset is \(n\)-expandable: every tested box of volume at least \(n\) contains at most \(k\leq n\) selected vertices, hence satisfies the expansion density bound. Set \[w'=\frac{n^{1-1/\alpha}}{2^d\beta}.\] Then \(s(w')=n\) and \(w'<w\). For all sufficiently large \(k\), \(w'\) exceeds Proposition 4.2’s fixed lower threshold. Apply that proposition to the selected subset at \(w'\). It gives a region of volume at least \(c k\). The actual sure-connectors have marks at least \(w\geq w'\), so monotonicity of the kernel preserves their connection probability lower bound. The same thinning calculation then gives \(\exp[-c'k w_{\rm hh}^{-(\tau-1)}]\), the intended heavy-heavy power. Increase the fixed minimum \(k\) to cover the auxiliary threshold, and treat the remaining finite values separately. The \(s(w)\leq n\) case remains the printed proof. This supplies the omitted case without claiming a false final merging result.\par
\section*{Independent checks and reference audit}
\subsection*{The coupling and the exponent calculation}
The sure/unsure construction in Definition 5.6 and Proposition 5.8, pp. 29–30, survives a direct distributional check. If \(s\) independent uniforms are split according to \(B=\{\min U_i>r\}\), with \(r=1-2^{-1/s}\), then \(\mathbb P(B)=1/2\). For \(0\leq z_i\leq1\), \[\mathbb P(U_i\leq z_i\ \forall i\mid B)
   =2\prod_i(z_i-r)_+,\qquad
 \mathbb P(U_i\leq z_i\ \forall i\mid B^c)
   =2\left(\prod_i z_i-\prod_i(z_i-r)_+\right).\] The equally weighted mixture is exactly \(\prod_i z_i\). Independent labels on the split Poisson process and independent edge vectors across connector vertices therefore give the required original graph law. Seven independently computed joint-distribution examples, evaluated with 70-digit working precision, agree with this identity. The algebraic identity, rather than numerical agreement, is the proof.\par
The high-mark scale satisfies \(k w_{\rm hh}^{-(\tau-1)}\asymp k^{\zeta_{\rm hh}}\). Balancing the expected removed-vertex count with the crossing-edge count gives \[\gamma_{\rm long}
  =\frac{\alpha-1}{\xi_\star+\tau-1},\qquad
 \xi_\star=\max\{0,\alpha-(\tau-1),
                  (\sigma+1)\alpha-2(\tau-1)\},
 \quad
 \gamma_\star=\min\{\gamma_{\rm long},1/(\sigma+1)\}.\] The identities and multiplicity cases of Lemma 7.6 and Appendix C were checked directly, including ties and the threshold limit. In particular, equality with the short-edge exponent introduces the additional logarithmic factor; it must not be discarded at a phase boundary. An independent exact-rational sweep checked 840 parameter cases and found no disagreement in the tested exponent and multiplicity identities. This finite sweep supplements the case algebra; it is not a universal proof by sampling.\par
The independent evidence file also contains lattice counts for five square boxes, the zero-edge and endpoint witnesses, the short/long-edge graph example, the Palm merge example, the derivative in Problem Problem 8 (\#prob:monotone), the cover-scale exponent \(11/7\), and the kernel expectation calculation below. The script was authored for this review and actually executed with exit code zero. No author code was executed. Exact expressions, parameters and execution evidence are retained under this paper’s verification tree.\par
\subsection*{Isolation, geometry and the external Erdős–Rényi input}
In the Poisson case, the mark-suppression regions are disjoint from the revealed low-mark rectangles once the size is sufficiently large. Their void probability is the exponential of their intensity. For the retained crossing edges, (7.21) bounds individual edge probabilities away from one; conditional edge independence and \(1-q\geq e^{-c q}\), followed by Jensen’s inequality, give (7.31)–(7.32), p. 48. This supplies the lower bound used in the main isolation argument. The lattice needs the additional nondegeneracy discussed in Problem Problem 1 (\#prob:lattice). The FKG step on p. 53 uses decreasing isolation and vertex-count events; the source first conditions on the point set and then applies the Poisson association step. Its sign is consistent with the required product lower bound.\par
The geometric discussion near (B.22), p. 67, informally describes an intersection of two spheres without spelling out a coarea Jacobian. A direct bound avoids relying on that description. At a fixed interior point and exterior boundary distance \(z_v\), the area of the exterior sphere inside a Euclidean ball of radius \(r\) about that point is at most \(C_d r^{d-1}\). Layer-cake integration of the kernel \(r^{-\alpha d}\), starting at distance \(z_u+z_v\), bounds its surface integral by \[C_{d,\alpha}(z_u+z_v)^{-d(\alpha-1)-1}.\] The integral converges because \(\alpha>1\). Combining this bound with the interior or exterior shell measure yields the radial powers used in (B.22)–(B.25). Thus the informal coordinate explanation is not treated here as evidence of a false isolation exponent. The constant-comparison repair for \(g\) remains necessary.\par
The proof of the backbone invokes O’Connell’s largest-component large-deviation theorem on p. 26. The acquired author-hosted reference [OConnell], physical pp. 1–5, includes Theorem 3.1 on printed p. 279. It concerns \({\rm ER}(m,c/m)\) at fixed \(c>0\); for \(c>1\), the rate function gives an exponential lower-tail estimate below the positive giant density. This matches the fixed-supercritical-mean use after the source’s domination step, with sufficiently large integer \(m\). Small-size conventions and rounding can be absorbed by increasing the minimum backbone scale. The whole nine-page reference proof was not independently certified.\par
\subsection*{Local convergence: finite mean is a separate hypothesis}
The law-of-large-numbers argument on pp. 40–41 cites Local limits of spatial inhomogeneous random graphs. The acquired version [LocalLimits] is arXiv:2107.08733v3 dated 6 September 2022, with a title-page date of 7 September; its first eight pages were read and viewed. Assumption 1.8(2), p. 5, and Theorems 1.9 and 1.11, pp. 6–7, require an integrable averaged connection bound. That hypothesis is not automatic for every \(\sigma\) allowed in the reviewed paper.\par
Equation reference: kernelmean. Indeed, with \(t=\tau-1>1\), integration over the ordered marks gives \[
 \mathbb E[\kappa_\sigma(W_1,W_2)]
  =\frac{2t^2}{t-1}\int_1^\infty u^{\sigma-2t}\,\mathrm du,\] which is finite exactly when \(\sigma<2\tau-3\). For finite \(\alpha>1\), spatial integration in [model] (\#model) is a positive constant times \(\beta p\,\alpha/(\alpha-1)\) times the kernel; for the threshold profile the corresponding factor is \(\beta p\). For example \(\tau=5/2,\sigma=3,\alpha=2\) is in the heavy-heavy regime, with \(\zeta_{\rm hh}=10/13>0\), but the integral in [kernelmean] (\#kernelmean) diverges. This shows that the checked reference theorem alone does not justify that full parameter range.\par
It does not show a failure of local convergence. For each fixed finite mark \(w\), \[\mathbb E[\kappa_\sigma(w,W)]
 =w\,\mathbb E[W^\sigma;W\leq w]
  +w^\sigma\,\mathbb E[W;W>w]<\infty.\] Thus the Palm degree is finite almost surely, and each fixed graph-radius neighborhood occupies a finite spatial region almost surely. This supplies a direct argument for the present Poisson restriction model. Fix a graph radius and a bounded rooted-neighborhood indicator. Approximate it by the same indicator in a Euclidean window of radius \(R\). The Palm probability that the full graph ball leaves this window tends to zero as \(R\to\infty\). For fixed \(R\), averages over well-separated spatial windows are independent, so their normalized empirical averages converge to the corresponding Palm expectation. Campbell’s formula bounds the average approximation error by the escape probability plus the vanishing boundary fraction of the large box. First let the ambient volume tend to infinity, then \(R\to\infty\). Finally \(N_n/n\to1\) converts volume normalization to uniform-vertex normalization. This proves the needed local convergence in probability without requiring a finite unconditional mean degree.\par
The random denominator in the uniform-vertex average on p. 41 should likewise be handled by \(N_n/n\to1\) before a Campbell/Palm comparison; conditioning on a uniformly selected point is not exactly the unweighted Palm law at finite volume. This concentration step supplies the asymptotic comparison needed there.\par
The acquired almost-local reference [AlmostLocal] is the later v3 of 6 October 2025. Its Theorem 2.2, pp. 4–5, requires local convergence in probability and a vanishing normalized count of pairs in distinct large components. The source’s outside-giant tail is sufficient for the pair condition, because at least one of such a pair is outside the largest component. The criterion also extends to random \(N_n\) with \(N_n/n\to1\): apply the same deterministic component-size inequalities with \(N_n\) in place of \(n\), and define the normalized quantities as zero on \(N_n=0\). The finite-mean/uniform-integrability assumption appearing in that reference’s edge-count result is not required for its largest-component result.\par
These checks record a reference-hypothesis qualification and supply the relevant argument for this model. They are not an allegation that Corollary 2.3 is false. The publisher PDF of the local-limits reference was not acquired, so no identity between its exact journal wording and the inspected preprint is claimed.\par
\section*{Version limits and overall assessment}
The frozen source manifest binds four files, with 2,883,644 aggregate bytes and no binding errors. Two PDF entries have the same main-source hash; they are not two independent versions. The arXiv HTML is a metadata record. The 1,056-byte response associated with the separate supplement DOI is an HTML security challenge, not a PDF. The embedded Appendices A–C were read in the 74-page source; equivalence with a separate journal supplement has not been established. A single ordinary attempt to retrieve an author-institution copy of the journal PDF failed with a connection timeout and produced no PDF. There was no bypass attempt. The findings therefore apply to the exact acquired preprint, without silently transferring them to an uninspected version of record.\par
A bounded primary-source correction search was performed. The current arXiv record identifies v3 with the same date. A 2026 publisher correction found for the related paper Large deviations of the giant in supercritical kernel-based spatial random graphs concerns editorial received/accepted dates and a different DOI. It is not a correction of the mathematical statements reviewed here. The later almost-local reference explicitly records a correction to its Corollary 2.6 and an added theorem; the actual acquired v3, rather than an assumed old statement, was used for this audit. No mathematical author correction to this exact assigned PDF was established by the bounded search. This is not a claim that no correction exists anywhere.\par
The detailed finite-\(n\) hyperbolic and fixed-number-of-vertices transfers mentioned on p. 54 are left partly to the reader in the source. Their complete model-dependent coupling and conditioning estimates were not reconstructed here. All external references beyond the specified acquired pages are treated as cited inputs, not as independently verified complete proofs. Missing supplements, unavailable publisher payloads and the absence of executed author code are verification limits, not mathematical defects.\par
Problem Problem 1 (\#prob:lattice) requires a correction to the stated range of a final theorem and to the associated lattice proof. The other numbered problems have explicit local repairs or restricted conclusions. In particular, the false event implication, incorrect conditioning field and invalid monotonicity explanation must be corrected even though their downstream asymptotic estimates can be recovered by the arguments given above. The checked exponent balances, conditional graph-law construction and Poisson isolation mechanism provide substantial counterevidence to a whole-paper failure claim. The scientific decision remains pending.\par
9 J. Jorritsma, J. Komjáthy and D. Mitsche, Cluster-size decay in supercritical kernel-based spatial random graphs, arXiv:2303.00724v3, 17 October 2024. Assigned journal DOI: 10.1214/24-AOP1742 (https://doi.org/10.1214/24-AOP1742). Exact preprint reviewed.\par
N. O’Connell, Some large deviation results for sparse random graphs, Probab. Theory Related Fields 110 (1998), 277–285. Acquired author-hosted PDF: https://maths.ucd.ie/\textasciitilde{}noconnell/pubs/noc98b.pdf. Physical pp. 1–5 inspected, including Theorem 3.1.\par
R. van der Hofstad, P. van der Hoorn and N. Maitra, Local limits of spatial inhomogeneous random graphs, arXiv:2107.08733v3, 6 September 2022. https://arxiv.org/abs/2107.08733v3. Physical pp. 1–8 inspected. Journal DOI: 10.1017/apr.2022.61 (https://doi.org/10.1017/apr.2022.61); journal bytes not matched.\par
R. van der Hofstad, The giant in random graphs is almost local, arXiv:2103.11733v3, 6 October 2025. https://arxiv.org/abs/2103.11733v3. Physical pp. 1–7 inspected, including Theorem 2.2 and the revised Corollary 2.6.\par
\end{document}
